\documentclass[11pt]{article}
\usepackage[margin=1.1in]{geometry}
\usepackage{amsmath,amssymb}
\usepackage{hyperref}
\usepackage{enumitem}
\hypersetup{colorlinks=true,linkcolor=blue,citecolor=blue,urlcolor=blue}

\newcommand{\refereequote}[1]{%
  \begin{quote}\itshape\small #1\end{quote}}

\begin{document}

\noindent
\textbf{Response to Referee Report (Second Referee)}\\
\textit{Relational Actualism and Quantum Mechanics (RAQM)}\\
Foundations of Physics

\bigskip\noindent
We thank the referee for a careful and technically precise report. The
major comments identify four specific gaps; we address each in turn,
followed by the minor comments.

\bigskip\hrule\bigskip

\noindent\textbf{Major comment A: S-matrix poles vs.\ localized DAG vertices}

\refereequote{There is a severe tension between claiming a discrete, localized
actualization event at a spacetime point $x_0^\mu$ and relying on an
$S$-matrix pole, which is a global property of the asymptotic future and
past. The precise local spacetime mechanism writing the DAG vertex remains
ambiguous.}

We have added a new paragraph to the localization section (\S\ref{sec:localization})
that resolves this tension directly.

The key point is that RA uses the on-shell condition not as a trigger
\emph{at} $t = \pm\infty$, but as a diagnostic for the \emph{type} of
interaction occurring at finite time. An interaction that would produce an
on-shell asymptotic state does so by creating an environmental entanglement
record at finite time $t_0$, well before $t = +\infty$.

The $S$-matrix pole characterises the \emph{amplitude} of the process; the
actualization event is the \emph{crossing} of the entropic threshold
$\Delta S > 0$ at finite time, triggered by the onset of asymptotically
irreversible environmental entanglement. The new paragraph includes a
concrete example:

\begin{quote}
``In a bubble chamber, the track forms at finite time and finite density ---
the on-shell condition is met locally when the produced particle propagates
far enough to generate an orthogonal environmental record, not at $t = \pm\infty$.
The LSZ asymptotic states are the mathematical idealization of this process;
the DAG vertex is the physical event at finite time.''
\end{quote}

The LSZ formula is the perturbative amplitude book-keeping tool; the
actualization criterion $\Delta S > 0$ is the physical trigger. These are
operating at different levels of description and do not conflict.

\bigskip\hrule\bigskip

\noindent\textbf{Major comment B: Finite-dimensional truncation and the
implicit cutoff}

\refereequote{The truncation scale or UV cutoff implicitly acts as a
fundamental parameter unless the Type III$_1$ formulation is fully
resolved. The truncation scale introduces an implicit free parameter.}

We have added a sentence directly following the $\mathcal{O}(1/d)$ error
estimate that closes this objection:

\begin{quote}
``Crucially, $d$ is not a free parameter of the theory: it is determined
by the physical size of the causal diamond $(T, R)$ of the experimental
setup. Every finite experiment specifies $(T, R)$ and thereby fixes $d$;
no truncation scale independent of the experiment is introduced.''
\end{quote}

The dimension $d$ of the truncated algebra is a property of the
\emph{experimental setup}, not a parameter of the RAQM framework.
The framework itself introduces no cutoff; the finite causal diamond of any
real measurement does. The Type III$_1$ extension (Hard Wall HW2) would
remove even this experiment-dependent truncation, promoting the approximate
argument to an exact theorem.

\bigskip\hrule\bigskip

\noindent\textbf{Major comment C: $\Delta S^*$ derivation relying on
unpublished companion papers}

\refereequote{This derivation relies entirely on external, unpublished
companion papers. Introducing specific numerical values like
$P_{\mathrm{acc}}(1) \approx 0.548$ reads as speculative and disjointed
from the core arguments. The paper would be stronger with a self-contained
derivation.}

The referee's concern is valid. We have added a self-contained description
of the $P_{\mathrm{acc}}(1)$ calculation directly in the revised text
(\S3.4):

\begin{quote}
``The key calculation: at density $\mu = 1$, the BDG-filtered Poisson-CSG
assigns to each vertex a score $S_{\mathrm{BDG}} = c_0 - c_1 N_1 + \cdots
+ c_4 N_4$ where $N_k \sim \mathrm{Poisson}(\mu^k/k!)$. The acceptance
probability $P_{\mathrm{acc}}(1) = \mathbb{P}[S_{\mathrm{BDG}} > 0]$ at
$\mu=1$ is evaluated by Monte Carlo simulation over the BDG integers
$(1,-1,9,-16,8)$, giving $P_{\mathrm{acc}}(1) \approx 0.548$. The Python
script is publicly available at the GitHub repository for independent
verification.''
\end{quote}

The computation is a straightforward Monte Carlo over the five BDG integers
with independent Poisson counts. The referee or any reader can reproduce it
in minutes from the public repository. The companion paper RASM provides
the theoretical context; the number itself is independently verifiable.

\bigskip\hrule\bigskip

\noindent\textbf{Major comment D: Joint snap probability for entangled pairs}

\refereequote{Since non-local actualization of entangled states is the primary
battleground of objective collapse models, relegating this to an open item
weakens the core thesis. At least a formal sketch is needed.}

We have replaced the previous "open item" note with a formal sketch
(\S5 open refinements, first item):

\begin{quote}
``A formal sketch of the resolution: the joint snap probability for a bipartite
entangled state $|\Psi_{AB}\rangle$ is obtained by applying the covariant rate
formula to the joint four-current $j^\mu_{AB}(x) = \langle \Psi_{AB}|
\hat{j}^\mu(x)|\Psi_{AB}\rangle$. By the amplitude\_locality axiom (each
local amplitude depends only on the causal past of the local region),
the joint snap rate factorises: $r_{AB}(x_A, x_B) = r_A(x_A) \cdot r_B(x_B)$
for spacelike-separated $x_A, x_B$. The joint snap probability is therefore
the product of the individual covariant rates, each of which is
frame-independent by the \texttt{frame\_independence} theorem. The joint
snap event is a spacelike-separated pair of independent actualization events,
consistent in every frame.''
\end{quote}

The factorisation relies on the amplitude\_locality axiom --- the same axiom
used in the causal invariance proof of RAEB. It is not derived here, but
it is the physically well-motivated assumption (QFT microcausality) that
grounds the no-signaling structure throughout.

\bigskip\hrule\bigskip

\noindent\textbf{Minor comment: Gauge invariance and dressed states}

\refereequote{Given that the primary bosons observed in nature are gauge
bosons, the manuscript should briefly discuss how physical Wilson lines or
dressed asymptotic states might resolve the edge-mode issue.}

We have added a sentence to the gauge invariance paragraph:

\begin{quote}
``A natural resolution in QED uses dressed asymptotic states
(Faddeev-Kulish coherent states), in which the charged particle is
accompanied by its soft-photon cloud; the corresponding Wilson-line
operator dresses the Fock-space state to produce an infrared-finite,
gauge-invariant asymptotic state on which $\Delta S > 0$ is well-defined.
Whether this extends cleanly to non-Abelian theories is an open target.''
\end{quote}

\noindent\textbf{Minor comment: Formal Poisson process definition}

\refereequote{A brief, formal mathematical definition of the competing
Poisson process would improve readability.}

A formal definition paragraph has been added immediately before the
covariant rate formula. It states the process precisely: each spacetime
volume element $d^4x$ fires independently at rate $r(x^\mu)$; the
probability of the first event occurring in $d^4x$ at $x^\mu$ is
$dP = r(x^\mu)\,d^4x \cdot \exp(-\int r\,d^4x')$; the ``competing''
character refers to the race in which the first volume element to fire
writes the DAG vertex.

\noindent\textbf{Minor comment: Companion paper dependence}

The self-contained $P_{\mathrm{acc}}$ calculation (Comment~C above),
the joint snap sketch (Comment~D), and the gauge invariance note all reduce
dependence on companion papers for load-bearing claims. The companion
citations remain for context and extended derivations, but the core
physical claims of this paper are now supported by material in the paper
itself.

\bigskip\hrule\bigskip

\noindent\textbf{Summary of changes (second referee)}

\begin{enumerate}[itemsep=4pt]

\item \textbf{New paragraph (\S\ref{sec:localization}): LSZ vs.\ finite-time
actualization.} Bubble chamber example; on-shell as finite-time diagnostic,
not asymptotic trigger.

\item \textbf{New sentence (state definitions)}: Truncation dimension $d$
is determined by $(T,R)$ of the experimental setup, not a free theory
parameter.

\item \textbf{Self-contained $P_{\mathrm{acc}}$ calculation (\S3.4)}:
BDG Poisson-CSG at $\mu=1$; Monte Carlo evaluation; GitHub verification link.

\item \textbf{Joint snap sketch (\S5)}: Amplitude\_locality $\Rightarrow$
factorisation of joint snap rate $\Rightarrow$ frame-independent joint
probability. Formal sketch replacing ``open item'' note.

\item \textbf{Gauge invariance note}: Faddeev-Kulish dressed states and
Wilson lines as the natural QED resolution; non-Abelian extension noted
as open.

\item \textbf{Competing Poisson process definition}: Formal paragraph
added before the covariant rate formula.

\end{enumerate}

\bigskip\noindent
We believe the revised manuscript now provides self-contained justification
for all load-bearing claims raised by this referee. The remaining
open items (Type III$_1$ extension, non-Abelian dressed states, amplitude
locality intrinsic proof) are precisely identified as Hard Walls with named
collaboration targets.

\bigskip
\noindent
\textit{Joshua F.\ Sandeman}\\
Independent Researcher, Salem, Oregon\\
\href{mailto:sansuikyo@gmail.com}{sansuikyo@gmail.com}

\end{document}
